\subsection{A universal two-constant action obstruction}

\begin{theorem}[Proved: the identity and two constants]
\label{CP2:A-two-constant-image}
Let $S$ be a nonempty finite left-$N$, right-$M$ biset, with
$M,N$ finite monoids. Suppose its right transformation image is
exactly
\[
\{\operatorname{id}_S,c_0,c_1\},
\]
where $c_i$ is constant at $s_i$ and $s_0\ne s_1$.
Then its two-object collage is not of type $\CP^2$.
The same conclusion holds with the two sides interchanged.
There is no bound on the cardinalities of $M,N,S$.
\end{theorem}
\begin{proof}
Commutation with the two constants forces every left actor to
fix $s_0$ and $s_1$. Put $U=\{s_0,s_1\}$ and let $X_U\subseteq X_S$
be the two collage spaces. Their corner spaces are identical.
The subset $U$ is a subbiset. Write
\[
\chi(m)=\begin{cases}1&\text{if $m$ acts identically on $S$},\\
0&\text{if $m$ acts constantly on $S$}.
\end{cases}
\]
This is a multiplicative zero-one character. On
$L=\ker(\mathbb Q[S]\to\mathbb Q)$ the right action is scalar
$\chi$. The vector $u=[s_1]-[s_0]$ spans a trivial left
$N$ submodule. Put $V=L/\mathbb Qu$ and let
\[
\gamma\in\operatorname{Ext}^1_{\mathbb QN}(V,\mathbb Q)
\]
be the class of
\begin{equation}\label{CP2:A-two-constant-extension}
0\longrightarrow\mathbb Q\xrightarrow{\,u\,}L
\longrightarrow V\longrightarrow0
\end{equation}
as left modules. In this display the commuting right action
on all three modules is scalar $\chi$.

Set
\[
D_N=\mathbb Q\mathbin{\otimes^{\mathbf L}_{\mathbb QN}}V,
\qquad
D_M=\mathbb Q_\chi\mathbin{\otimes^{\mathbf L}_{\mathbb QM}}\mathbb Q.
\]
The relative bridge coefficient is
$\mathbb Q[S]/\mathbb Q[U]=V\otimes\mathbb Q_\chi$.
The double mapping-cylinder chains therefore identify the
reduced relative complex of $(X_S,X_U)$ with
$(D_N\otimes_{\mathbb Q}D_M)[1]$.
Since $\chi$ takes the value zero, $H_0(D_M)=0$.
In particular K\"unneth supplies the summand
\[
\operatorname{Tor}^{\mathbb QN}_1(\mathbb Q,V)
\otimes H_1(D_M)
\ \subseteq\ H_3(X_S,X_U;\mathbb Q).
\]

We identify its boundary, so no comparison of dimensions alone
is being used. The coefficient extension for $U\subset S$ is
the pushout of \eqref{CP2:A-two-constant-extension}, tensored
on the right by $\mathbb Q_\chi$, along the inclusion
$u:\mathbb Q_\chi\to\mathbb Q[U]$.
Choose any vector-space splitting of $L\to V$. In the resulting
normalized double-bar chains the off-diagonal left-action face
is exactly a cocycle for $\gamma$, the right-action faces are
scalar $\chi$, and the vertex augmentations vanish on these
augmentation-module lifts. Hence the connecting chain map is
$\gamma\otimes\operatorname{id}$ followed by the actual
inclusion of $u$ in the core complex.

The left $N$ action on $U$ is trivial. By
Lemma~\ref{CP2:A-trivial-biset-side-H2}, collapsing that corner
induces an isomorphism on rational second homology. The
augmentation sequence
\[
0\longrightarrow\mathbb Q_\chi\xrightarrow{\,u\,}
\mathbb Q[U]\longrightarrow\mathbb Q\longrightarrow0
\]
identifies the reduced ordinary core act-cone complex with
$D_M[1]$, naturally with respect to the displayed inclusion.
Consequently the relative boundary on the preceding K\"unneth
summand is, up to the harmless global boundary sign,
\begin{equation}\label{CP2:A-two-constant-boundary}
\gamma\otimes\operatorname{id}:
\operatorname{Tor}^{\mathbb QN}_1(\mathbb Q,V)
\otimes H_2(X_U;\mathbb Q)
\longrightarrow H_2(X_U;\mathbb Q).
\end{equation}
The duality being used here is the exact vector-space duality
\[
\operatorname{Ext}^1_{\mathbb QN}(V,\mathbb Q)
\cong
\operatorname{Hom}_{\mathbb Q}
\bigl(\operatorname{Tor}^{\mathbb QN}_1(\mathbb Q,V),\mathbb Q\bigr).
\]
It follows by applying $\operatorname{Hom}_{\mathbb Q}(-,\mathbb Q)$
to a free resolution of $V$ tensored with the trivial right
module; the resulting cochain complex is
$\operatorname{Hom}_{\mathbb QN}(-,\mathbb Q)$.

If $\gamma\ne0$, it is a nonzero, hence surjective, functional.
Equation~\eqref{CP2:A-two-constant-boundary} is therefore onto.
The map $H_2(X_U;\mathbb Q)\to H_2(X_S;\mathbb Q)$ vanishes,
so every rational degree-two class on $X_S$ is relative to $X_U$.
Both corners are contained in $X_U$; collapsing it in the
mapping cylinder gives the reduced suspension of the quotient
of the two edge-bar spaces. All positive-degree products in
that quotient vanish. Their pullbacks to $X_S$ therefore vanish,
which is incompatible with the square on $\CP^2$.

If $\gamma=0$, take an $N$-linear retraction
$q:L\to\mathbb Q$ with $q(u)=1$ and put
$q(s)=q([s]-[s_0])$. It satisfies
$q(s_0)=0$, $q(s_1)=1$ and $q(ns)=q(s)$.
Let $\rho:\mathbb Q\to U$ send $1$ to $s_1$ and every other
rational number to $s_0$. Then $r(s)=\rho(q(s))$ fixes $U$
and is left equivariant. It is also right equivariant, since
every right actor is identity or constant at one of the two
fixed points. Thus $r:S\to U$ is an actual biset retraction,
and $X_U$ is an actual retract of $X_S$.

Every rational degree-two product on $X_U$ is zero: its
trivial-side collapse is a degree-two cohomology isomorphism
to the ordinary two-state $M$-act cone, and that act embeds
in the affine line as $\{0,1\}$, with identity and constant
affine maps. Apply Lemma~\ref{CP2:A-high-affine} and pull back.
If $X_S$ were a projective plane, the split injection
$H^2(X_U;\mathbb Q)\to H^2(X_S;\mathbb Q)$ would have rank
zero or one. In the zero case its degree-two generator is
relative to $X_U$ and has square zero by the suspension
quotient. In the rank-one case the generator is pulled back
from $X_U$ through the retraction and again has square zero.
Both are contradictions. This proves the theorem; opposition
proves its left-handed version.
\end{proof}

\begin{remark*}[Computational scope of this gate]
Testing the literal transformation image against exactly the
three maps in the theorem is sufficient. No corner monoid is
replaced by that image, and no homotopy-type assumption on
an individual corner is made. The split case asserts the
projective-plane obstruction; it does not assert that all
products on an arbitrary larger-rank collage vanish.
\end{remark*}
